% Spanish summary, included for completeness even though the report is
% written in English. Mirrors the English abstract.
\renewcommand{\abstractname}{Resumen}
\begin{abstract}
% TODO redactar y ajustar a 150-300 palabras.
Este Trabajo Fin de Grado presenta una prueba de concepto de un
orquestador universal de máquinas virtuales, orientado a entornos
docentes y de investigación. El sistema proporciona una capa de
abstracción común que permite gestionar máquinas virtuales con
independencia del hipervisor subyacente, integrando en un único modelo
un árbol de descriptores con herencia en cascada, un motor de
operaciones basado en el patrón Command y una API REST como frontera
pública. La arquitectura se diseña en dos capas: un núcleo genérico
reutilizable y una aplicación específica para despliegue de
asignaturas, materializada como capa cliente sobre el primero.

El sistema se entrega como un contenedor Docker descargable que el
administrador despliega con un instalador de una sola línea, junto con
un cliente de línea de órdenes que opera tanto en modo interactivo
como por instrucciones sueltas. La adaptación a cada hipervisor se
realiza a través de conectores; en esta versión se incluye un
conector simulado que permite ejercitar el modelo completo sin
infraestructura real, sirviendo como punto de partida para los
conectores reales contra VMware y Proxmox planteados como continuación
del trabajo.

% TODO confirmar palabras clave finales.
\medskip

\noindent\textbf{Palabras clave:} orquestación de máquinas virtuales,
infraestructura como código, entornos docentes, herencia en cascada,
arquitectura por capas.
\end{abstract}
